%% kapa: exam=1 round=27 year=2016 period=2 form=A subject=appraisal_law subject_ko=감정평가관계법규 seq=17 number=17
%% source: Q-Net 감정평가사 기출문제 (cst00309) artlSeq=5169402 file=extracted/1008/5169402/제27회 감정평가사 1차 문제/제27회 감정평가사 2교시문제(A형).pdf
%% conversion: kordoc 4.24.51 -> tools/gen1.py (auto); review=proofread-sonnet; flags=none
\begin{question}{17}
부동산가격공시 및 감정평가에 관한 법령상 감정평가법인에 관한 설명으로 옳은 것을 모두 고른 것은?
\begin{framedtext}
ㄱ. 감정평가법인에는 5명 이상의 감정평가사를 두어야 한다. \par
ㄴ. 감정평가법인의 주사무소에 주재하는 최소 감정평가사의 수는 3명이고, 분사무소에 주재하는 최소 감정평가사의 수는 2명이다. \par
ㄷ. 감정평가법인이 해산하고자 할 때에는 국토교통부장관의 인가를 받아야 한다. \par
ㄹ. 감정평가법인에 대해서는「부동산가격공시 및 감정평가에 관한 법률」에서 정한 것을 제외하고는「상법」중 회사에 관한 규정을 준용 한다. \par
\end{framedtext}
\begin{choices}
  \item ㄱ, ㄴ
  \item ㄱ, ㄷ
  \item ㄱ, ㄹ
  \item ㄴ, ㄷ
  \item ㄴ, ㄹ
\end{choices}
\end{question}
\answer{3}
